% GENERATED FILE. Edit canonical lessons, then run npm run generate:books.
% canonical-source-hash: fnv1a64:354f8e8bfd02f044
% canonical-lessons: ZH-C180-zuoke, ZH-C180-zhuren, ZH-C180-kanwang, ZH-C180-guolai, ZH-C180-yue

\chapter{Coming Over}
\label{ch:zh-coming-over}
\hlchaptermodality{180}

\begin{chapteropening}
\textbf{By the end of this chapter:} \emph{I can talk about visiting as a guest, the host, calling on someone, coming over and making a date.}

Complete the last lesson of chapter 180: I can talk about visiting as a guest, the host, calling on someone, coming over and making a date.
\end{chapteropening}

\section[\texorpdfstring{\zh{做客} (zuòkè)}{做客 (zuòkè)}]{\zh{做客} (zuòkè) — to visit as a guest; to be a guest}

\label{lesson:ZH-C180-zuoke}



\begin{quote}
Before the new one: say the Mandarin for to iron, then the Mandarin for a wardrobe.
\end{quote}

\subsection*{You'll want to know: \zh{做客}}
\textbf{\zh{做客}} — \emph{zuòkè} — \textquotedblleft{}to visit as a guest; to be a guest\textquotedblright{}.

\begin{quote}
\textbf{\zh{我去他家做客喝了两杯茶}} \emph{wǒ qù tā jiā zuòkè hē le liǎng bēi chá} — I went round to his place and had two cups of tea.
\end{quote}

\begin{cousinweb}
\textbf{\zh{做客}} (\emph{zuòkè}) is built from \zh{做} 'to do, to be' and \zh{客} 'guest'. The second is inherited from Middle Chinese \emph{khaek} (Guangyun, in Baxter's transcription), glossed 'guest'; \zh{做} has no form in that dictionary, so the history of the first part is not settled.
\end{cousinweb}

\begin{grammarlens}[title={What you've built}]
One more word for talking about plans.
\end{grammarlens}

\subsection*{Your turn}
\begin{itemize}
\raggedright
  \item \emph{Say it:} \emph{zuòkè}
  \item \emph{Say it:} \emph{zuòkè}, once more
  \item \emph{Say it:} \emph{wǒ qù tā jiā zuòkè hē le liǎng bēi chá} — I went round to his place and had two cups of tea.
  \item \emph{Recall:} say the Mandarin for to iron, then the Mandarin for a wardrobe, then say \emph{zuòkè} again
\end{itemize}

\begin{tcolorbox}[breakable,colback=teal!4,colframe=teal!35!black,title={Before you move on}]
What does \zh{做客} mean? (\textquotedblleft{}To visit as a guest; to be a guest\textquotedblright{}.) Say it once more.
\end{tcolorbox}
\section[zhǔrén]{zhǔrén — host; owner}

\label{lesson:ZH-C180-zhuren}



\begin{quote}
Before the new one: say the Mandarin for a wardrobe, then the Mandarin for to visit as a guest.
\end{quote}

\subsection*{You'll want to know: zhǔrén}
\emph{zhǔrén} — \textquotedblleft{}host; owner\textquotedblright{}.

\begin{quote}
\emph{zhǔrén hěn rèqíng} — The host was very welcoming.
\end{quote}

Its characters come later; for now it is a word for the ear.

\begin{cousinweb}
\emph{zhǔrén} is built from \emph{zhǔ} 'to be in charge' and \zh{人} 'person', both inherited from Middle Chinese, \emph{tsyuX} and \emph{nyin} (Guangyun, in Baxter's transcription): the person in charge of the house.
\end{cousinweb}

\begin{grammarlens}[title={What you've built}]
One more word for talking about plans.
\end{grammarlens}

\subsection*{Your turn}
\begin{itemize}
\raggedright
  \item \emph{Say it:} \emph{zhǔrén}
  \item \emph{Say it:} \emph{zhǔrén}, once more
  \item \emph{Say it:} \emph{zhǔrén hěn rèqíng} — The host was very welcoming.
  \item \emph{Recall:} say the Mandarin for a wardrobe, then the Mandarin for to visit as a guest, then say \emph{zhǔrén} again
\end{itemize}

\begin{tcolorbox}[breakable,colback=teal!4,colframe=teal!35!black,title={Before you move on}]
What does zhǔrén mean? (\textquotedblleft{}Host; owner\textquotedblright{}.) Say it once more.
\end{tcolorbox}
\section[kànwàng]{kànwàng — to call on; to visit (a person)}

\label{lesson:ZH-C180-kanwang}



\begin{quote}
Before the new one: say the Mandarin for to visit as a guest, then the Mandarin for the host.
\end{quote}

\subsection*{You'll want to know: kànwàng}
\emph{kànwàng} — \textquotedblleft{}to call on; to visit (a person)\textquotedblright{}.

\begin{quote}
\emph{zhōumò wǒ qù kànwàng wǒ de lǎoshī} — At the weekend I'm going to call on my teacher.
\end{quote}

Its characters come later; for now it is a word for the ear.

\begin{cousinweb}
\emph{kànwàng} is built from \zh{看} 'to look' and \emph{wàng} 'to gaze', both inherited from Middle Chinese, \emph{khanH} and \emph{mjangH} (Guangyun, in Baxter's transcription), which glosses the second with this very word.
\end{cousinweb}

\begin{grammarlens}[title={What you've built}]
One more word for talking about plans.
\end{grammarlens}

\subsection*{Your turn}
\begin{itemize}
\raggedright
  \item \emph{Say it:} \emph{kànwàng}
  \item \emph{Say it:} \emph{kànwàng}, once more
  \item \emph{Say it:} \emph{zhōumò wǒ qù kànwàng wǒ de lǎoshī} — At the weekend I'm going to call on my teacher.
  \item \emph{Recall:} say the Mandarin for to visit as a guest, then the Mandarin for the host, then say \emph{kànwàng} again
\end{itemize}

\begin{tcolorbox}[breakable,colback=teal!4,colframe=teal!35!black,title={Before you move on}]
What does kànwàng mean? (\textquotedblleft{}To call on; to visit (a person)\textquotedblright{}.) Say it once more.
\end{tcolorbox}
\section[\texorpdfstring{\zh{过来} (guòlai)}{过来 (guòlai)}]{\zh{过来} (guòlai) — to come over}

\label{lesson:ZH-C180-guolai}



\begin{quote}
Before the new one: say the Mandarin for the host, then the Mandarin for to call on someone.
\end{quote}

\subsection*{You'll want to know: \zh{过来}}
\textbf{\zh{过来}} — \emph{guòlai} — \textquotedblleft{}to come over\textquotedblright{}.

\begin{quote}
\textbf{\zh{你什么时候过来}} \emph{nǐ shénme shíhou guòlai} — When are you coming over?
\end{quote}

\begin{cousinweb}
\textbf{\zh{过来}} (\emph{guòlai}) is built from \zh{过} 'to pass' and \zh{来} 'to come', both inherited from Middle Chinese, \emph{kwaH} and \emph{loj} (Guangyun, in Baxter's transcription): to pass over to here. In the word the second loses its tone.
\end{cousinweb}

\begin{grammarlens}[title={What you've built}]
One more word for talking about plans.
\end{grammarlens}

\subsection*{Your turn}
\begin{itemize}
\raggedright
  \item \emph{Say it:} \emph{guòlai}
  \item \emph{Say it:} \emph{guòlai}, once more
  \item \emph{Say it:} \emph{nǐ shénme shíhou guòlai} — When are you coming over?
  \item \emph{Recall:} say the Mandarin for the host, then the Mandarin for to call on someone, then say \emph{guòlai} again
\end{itemize}

\begin{tcolorbox}[breakable,colback=teal!4,colframe=teal!35!black,title={Before you move on}]
What does \zh{过来} mean? (\textquotedblleft{}To come over\textquotedblright{}.) Say it once more.
\end{tcolorbox}
\section[yuē]{yuē — to arrange to meet; to make a date with}

\label{lesson:ZH-C180-yue}



\begin{quote}
Before the new one: say the Mandarin for to call on someone, then the Mandarin for to come over.
\end{quote}

\subsection*{You'll want to know: yuē}
\emph{yuē} — \textquotedblleft{}to arrange to meet; to make a date with\textquotedblright{}.

\begin{quote}
\emph{wǒ yuē le péngyou xiàwǔ jiànmiàn} — I've arranged to meet a friend this afternoon.
\end{quote}

Its characters come later; for now it is a word for the ear.

\begin{cousinweb}
\emph{yuē} is inherited from Middle Chinese \emph{'jak} (Guangyun, in Baxter's transcription), glossed 'to bind' and 'sparing': to bind someone to a time.
\end{cousinweb}

\begin{grammarlens}[title={What you've built}]
One more word for talking about plans.
\end{grammarlens}

\subsection*{Your turn}
\begin{itemize}
\raggedright
  \item \emph{Say it:} \emph{yuē}
  \item \emph{Say it:} \emph{yuē}, once more
  \item \emph{Say it:} \emph{wǒ yuē le péngyou xiàwǔ jiànmiàn} — I've arranged to meet a friend this afternoon.
  \item \emph{Recall:} say the Mandarin for to call on someone, then the Mandarin for to come over, then say \emph{yuē} again
\end{itemize}

\begin{tcolorbox}[breakable,colback=teal!4,colframe=teal!35!black,title={Before you move on}]
What does yuē mean? (\textquotedblleft{}To arrange to meet; to make a date with\textquotedblright{}.) Say it once more.
\end{tcolorbox}
